\section{Introduction}
\label{sec:intro}

Professional traders do not make decisions by scanning columns of numbers. They read charts. A practiced eye can identify a head-and-shoulders formation, a trendline break, or a volume divergence in seconds---patterns that encode complex multi-scale temporal and structural information in a spatial representation that humans process efficiently. Technical chart analysis has been a cornerstone of trading practice for over a century, with well-documented chart patterns serving as the shared visual language of financial markets~\cite{murphy1999technical}.

Despite the centrality of visual chart reasoning in human financial decision-making, the dominant paradigm in AI-driven financial prediction treats price data exclusively as numerical time series. State-of-the-art sequence models---from LSTMs and Transformers to recent architectures such as PatchTST~\cite{nie2023patchtst} and TimesNet~\cite{wu2023timesnet}---process raw OHLCV vectors directly, discarding the spatial-relational structure that emerges when this data is rendered as a chart. Meanwhile, the recent surge of interest in vision-language models (VLMs) and LLM-based financial agents has produced systems that either analyze text-based market data~\cite{fincon2024,tradingagents2024} or, at best, feed raw chart screenshots to generic visual models without any structured understanding of chart geometry.

We argue that this gap represents a missed opportunity. Chart visualizations are not merely cosmetic representations of numerical data---they are structured spatial encodings that make certain patterns perceptually salient. A double-bottom formation, for instance, is trivially visible in a candlestick chart but requires complex feature engineering to detect from raw price sequences. Support and resistance levels emerge as visually obvious horizontal bands but correspond to subtle distributional properties of price history. The insight that ``a picture is worth a thousand time steps'' has been underexploited in financial AI.

In this paper, we introduce \textsc{FinVL-MAS} (Financial Visual-Linguistic Multi-Agent System), a framework that operationalizes the cognitive workflow of expert traders as a multi-agent system for financial decision-making. Rather than treating chart images as opaque inputs to a generic VLM, \textsc{FinVL-MAS} extracts \emph{structured chart geometry}---trendlines, support/resistance levels, candlestick formations, chart patterns, and volume signals---and routes these structured representations through a team of cognitively-specialized reasoning agents coordinated via gated shared memory.

The multi-agent decomposition is not merely an engineering convenience; it reflects a testable hypothesis about the structure of expert trader cognition. We decompose the decision process into five functional stages---\emph{perception} (chart geometry extraction), \emph{pattern recognition} (formation matching and interpretation), \emph{contextual integration} (event and fundamental analysis), \emph{risk assessment} (position sizing and stop-loss determination), and \emph{decision formation} (evidence integration and action selection)---each implemented as a specialist agent. This decomposition enables systematic ablation of each reasoning component and produces interpretable decision rationales.

\paragraph{Contributions.} We make the following contributions:
\begin{itemize}
    \item We propose a \textbf{structured chart geometry extraction} pipeline that combines rule-based geometric detection with VLM-based holistic pattern recognition, producing typed representations of trendlines, support/resistance levels, candlestick patterns, and chart formations.
    \item We design a \textbf{cognitively-motivated multi-agent architecture} with five specialist agents that decompose the expert trader's reasoning process, coordinated via gated cross-attention shared memory.
    \item We provide \textbf{rigorous ablation evidence} isolating the marginal value of visual chart reasoning over numerical baselines, with regime-conditioned analysis showing when visual reasoning is most valuable.
    \item We adapt the \textbf{closed-loop R\&D workflow} from R\&D-Agent~\cite{rdagent2025} for iterative hypothesis testing about which visual patterns are most predictive under specific market conditions.
    \item We evaluate under \textbf{realistic financial constraints}---transaction costs, temporal splits, point-in-time alignment, and execution delays---and show that structured visual reasoning produces better risk-adjusted returns than state-of-the-art baselines, particularly during volatile market regimes.
\end{itemize}
